\chapter{From an Idea to a Validated Simulation}
\label{chap:overview}

The internship implementation appears only as a running example. The workflow
itself applies to behavior cloning, action-chunking models, diffusion-based
policies, or another vision-conditioned controller.

Use \filepath{README_ENV.md}, \filepath{README_SERVER.md}, and
\filepath{README_CODE.md} for installation and copyable project commands; this
guide focuses on the decisions and checks that remain useful across systems.

\section{The experimental cycle}

Most robot-learning projects can be organized as the following loop:

\begin{center}
\resizebox{\textwidth}{!}{%
\begin{tikzpicture}[
  node distance=8mm and 7mm,
  box/.style={draw=lirmmblue,rounded corners=1mm,thick,align=center,
              minimum height=10mm,minimum width=26mm,fill=softblue},
  arr/.style={-{Latex[length=2.2mm]},thick,lirmmblue}
]
  \node[box] (task) {Define a\\benchmark};
  \node[box,right=of task] (sim) {Build and test\\in simulation};
  \node[box,right=of sim] (data) {Collect and\\inspect data};
  \node[box,right=of data] (train) {Train and select\\a checkpoint};
  \node[box,right=of train] (deploy) {Deploy and\\evaluate};
  \draw[arr] (task) -- (sim);
  \draw[arr] (sim) -- (data);
  \draw[arr] (data) -- (train);
  \draw[arr] (train) -- (deploy);
  \draw[arr] (deploy.south) |- ++(0,-8mm) -| (task.south);
\end{tikzpicture}%
}
\end{center}

Each stage should produce evidence before the next one starts. A scene that
loads is not yet a valid benchmark; a dataset that has the expected shape is
not necessarily useful; a low validation loss is not proof of closed-loop task
success; and a robot that moves is not yet a safe deployment.

\begin{methodbox}[General rule]
Change one source of difficulty at a time. Begin with a task that is easy to
interpret, then add position, orientation, lighting, object, or geometry
variation progressively. When several variables change together, a failure is
hard to diagnose and the next dataset is difficult to design.
\end{methodbox}

\section{Define success before writing code}

A benchmark description should answer five questions:

\begin{enumerate}
  \item What is the initial state distribution?
  \item What observations are available to the policy?
  \item What actions can the policy command?
  \item What event counts as success or failure?
  \item Which conditions remain fixed during one evaluation?
\end{enumerate}

Avoid descriptions such as ``the robot should usually grasp the object.'' A
useful definition is observable: for example, ``the object starts inside a
marked region and the trial succeeds when it is released inside the target box
without triggering a safety stop.'' Also define a maximum duration and the
handling of contacts, re-grasps, human intervention, and sensor failure.

\begin{projectbox}[Internship example]
The project moved from simulation to a simple physical pick-and-drop task, then
to constrained extraction with progressively larger pose variation. Each stage
added one observable difficulty while retaining the same experimental cycle.
\end{projectbox}

\section{Separate interfaces from implementations}

A learned policy should be treated as one component with a clear contract:

\begin{equation}
  \text{recent observations} \longrightarrow
  \text{future action sequence}.
\end{equation}

The simulator and physical robot do not need the same internal code, but they
must provide compatible concepts: camera images, robot state, action decoding,
and timing. Keeping these interfaces explicit makes it possible to replace the
model architecture or manipulator without rewriting the complete workflow.

Simulation can validate the method and interfaces without implying direct
transfer of one trained model to the physical robot. Decide explicitly whether
simulation data, model weights, or only software interfaces will be reused.

\section{What should remain reproducible}

For every dataset, training run, and evaluation, record at least:

\begin{itemize}
  \item exact code version (the result of \filepath{git rev-parse HEAD}) and
  configuration values used for the run;
  \item task definition and object-pose range;
  \item camera placement, crop, resolution, and exposure mode;
  \item robot start state and control mode;
  \item observation and action definitions;
  \item number of accepted demonstrations;
  \item selected checkpoint and deployment parameters;
  \item individual evaluation trials, not only the final percentage.
\end{itemize}

The exact command is part of the experiment. A checkpoint filename alone is
insufficient because horizons, camera preprocessing, action convention, and
safety limits may be supplied at runtime.

\begin{checkpointbox}[Before continuing]
Write a one-page benchmark protocol and draw the observation/action interface.
If success, variation, and action semantics are still ambiguous, do not begin a
large data collection.
\end{checkpointbox}
